% Must be set before xcolor is first loaded (by beamer core) to avoid an
% "Option clash for package xcolor" error later when we request [table].
\PassOptionsToPackage{table}{xcolor}

\documentclass[9pt,aspectratio=169]{beamer}
\setbeameroption{show notes on second screen=right}
%% Choose aspect ratio and other standard options:
% [aspectratio=169] % 16:9 (default)
% [aspectratio=43]  % 4:3

\usetheme[minimal]{tugraz2018}
%% Choose main theme variant:
% [standard]        % standard (default)
% [institute]       % with institute's graphical acronym on the left
% [minimal]         % with reduced visuals

\usepackage[utf8]{inputenc}
\usepackage[english]{babel}

\usepackage[table]{xcolor}
\usepackage{booktabs}
\usepackage{tikz}
\usetikzlibrary{arrows.meta, positioning}
\usepackage{amssymb} % \checkmark
\usepackage{textcomp} % \texttimes

% ---------------------------------------------------------------------
% Our own images (Formations/, results/) live in the main project
% folder, three levels up from this template folder.
% ---------------------------------------------------------------------
\newcommand{\projroot}{../../../}

% ---------------------------------------------------------------------
% Marks any value that depends on the (currently still running) model
% evaluation. Search for \PENDING to find every spot to fill in once the
% notebook finishes. Takes no argument, so always follow it either with
% a closing parenthesis/punctuation, or with {} before more running text
% (bare control words swallow the space after them).
% ---------------------------------------------------------------------
\newcommand{\PENDING}{\textcolor{tug}{\textbf{[TBD]}}}

%% Enter presentation metadata
\title[Formation Classification]{Zero-Shot Formation Classification in Amateur American Football}
\author{Lukas Brandner}
\date{08.10.26}
\institute{IML}
\instituteurl{www.tugraz.at/institute/iml}

\begin{document}

\begin{frame}[plain]
  \maketitle
\end{frame}

\begin{frame}{Outline}
  \tableofcontents
\end{frame}

% =====================================================================
\section{Motivation}
% =====================================================================

\begin{frame}{Why automate formation tagging?}
  \begin{itemize}
    \item American football is ``chess on the field'': pre-snap formations carry tactical information coaches rely on
    \item At the amateur level (e.g.\ Austrian Football League), formation tagging on platforms like Hudl is done \textbf{entirely by hand}
    \item Open-weight Vision-Language Models (VLMs) can be run locally, could they help?
  \end{itemize}
\end{frame}

\begin{frame}{Research Questions}
  \begin{itemize}
    \item[\textbf{RQ1}] To what extent can user-grade open-source VLMs identify offensive formations from pre-snap images in a zero-shot setting?
    \item[\textbf{RQ2}] How do Qwen-VL, Gemma, and LLaMA compare in their ability to detect formations and player alignments?
    \item[\textbf{RQ3}] What are the main failure modes and limitations when applying VLMs to amateur-quality football broadcast footage?
  \end{itemize}
  \note{
  	\begin{itemize}
  		\item RQ1: can user-grade open-source VLMs identify formations zero-shot?
  		\item RQ2: how do the three models (Qwen-VL, Gemma, LLaMA) compare?
  		\item RQ3: what are the main failure modes on amateur-quality footage ?
  	\end{itemize}
  }
\end{frame}

% =====================================================================
\section{Background}
% =====================================================================

\begin{frame}{What a pre-snap formation shows}
  \begin{columns}[T]
    \begin{column}{0.5\textwidth}
      \small
      \begin{itemize}
        \item Backfield occupied (RB/FB), or empty?
        \item How many wide receivers per side?
      \end{itemize}
    \end{column}
    \begin{column}{0.5\textwidth}
      \small
      \begin{itemize}
        \item Tight end/fullback on the line?
        \item Which side has the extra receiver (``strong side'')?
      \end{itemize}
    \end{column}
  \end{columns}
  \vfill
  {\small\textbf{Why this is hard:} precise counting and spatial reasoning}
  \note{
 		- Four visual cues define a formation: backfield occupied or empty, WR count
 		per side, TE/FB on the line or not, which side has the extra receiver
 		("strong side").\\
 		- Why it's hard: needs precise counting + spatial reasoning — easy to miss
 		even for a person on grainy broadcast footage.
 		}
\end{frame}

\begin{frame}{The Formation Classes (1/2)}
  \centering
  \begin{columns}[T,onlytextwidth]
    \begin{column}{0.33\textwidth}
      \centering
      \includegraphics[width=0.95\linewidth]{\projroot Formations/EMPTY.png}\\
      \textbf{EMPTY}
    \end{column}
    \begin{column}{0.33\textwidth}
      \centering
      \includegraphics[width=0.95\linewidth]{\projroot Formations/TRIO.png}\\
      \textbf{TRIO}
    \end{column}
    \begin{column}{0.33\textwidth}
      \centering
      \includegraphics[width=0.95\linewidth]{\projroot Formations/DOUBLES.png}\\
      \textbf{DOUBLES}
    \end{column}
  \end{columns}
\end{frame}

\begin{frame}{The Formation Classes (2/2)}
  \centering
  \begin{columns}[T,onlytextwidth]
    \begin{column}{0.1\textwidth}\end{column}
    \begin{column}{0.35\textwidth}
      \centering
      \includegraphics[width=0.95\linewidth]{\projroot Formations/TREY-FB.png}\\
      \textbf{TREY}
    \end{column}
    \begin{column}{0.35\textwidth}
      \centering
      \includegraphics[width=0.95\linewidth]{\projroot Formations/TANK-FB.png}\\
      \textbf{TANK}
    \end{column}
    \begin{column}{0.1\textwidth}\end{column}
  \end{columns}
\end{frame}

\begin{frame}{Four Games, Four Different Cameras: Up Close}
  \centering
  \begin{columns}[T,onlytextwidth]
    \begin{column}{0.48\textwidth}
      \centering
      \includegraphics[height=4.1cm]{\projroot Testset/Presnap_Images/Game1_Wide---Clip-008_t12.jpg}\\
      \textbf{Game 1}
    \end{column}
    \begin{column}{0.48\textwidth}
    	\centering
    	\includegraphics[height=4.1cm]{\projroot Testset/Presnap_Images/Game3_Wide---Clip-006_t8.jpg}\\
    	\textbf{Game 3}
    \end{column}
  \end{columns}
\end{frame}

\begin{frame}{Four Games, Four Different Cameras: From Further Away}
  \centering
  \begin{columns}[T,onlytextwidth]
    \begin{column}{0.48\textwidth}
      \centering
      \includegraphics[height=4.1cm]{\projroot Testset/Presnap_Images/Game2_Wide---Clip-003_t2.jpg}\\
      \textbf{Game 2}
    \end{column}
    \begin{column}{0.48\textwidth}
      \centering
      \includegraphics[height=4.1cm]{\projroot Testset/Presnap_Images/Game4_Wide---Clip-020_t16.jpg}\\
      \textbf{Game 4}
    \end{column}
  \end{columns}
\end{frame}


% =====================================================================
\section{Data \& Methods}
% =====================================================================

\begin{frame}{Dataset}
  \begin{columns}[T]
    \begin{column}{0.55\textwidth}
      \small
      \begin{itemize}
        \item 40 offensive plays, 4 AFL games
        \item One pre-snap frame per play
        \item Manual annotations: formation, play type, jersey colors, offense side
        \item Majority class: \textbf{TRIO} (35\%)
      \end{itemize}
    \end{column}
    \begin{column}{0.45\textwidth}
      \centering
      \small
      \renewcommand{\arraystretch}{1.0}
      \begin{tabular}{lrr}
        \toprule
        \textbf{Formation} & \textbf{Count} & \textbf{\%} \\
        \midrule
        TRIO    & 14 & 35.0 \\
        DOUBLES & 11 & 27.5 \\
        EMPTY   &  5 & 12.5 \\
        TREY    &  5 & 12.5 \\
        TANK    &  5 & 12.5 \\
        \midrule
        Total   & 40 & 100.0 \\
        \bottomrule
      \end{tabular}
    \end{column}
  \end{columns}
  \note{- 40 offensive plays across 4 AFL games, one pre-snap frame per play.\\
  	- Manually annotated: formation, play type, jersey colors, offense side.\\
  	- Majority class TRIO at 35\%}
\end{frame}

\begin{frame}{Models Under Test}
  \centering
  \renewcommand{\arraystretch}{1.3}
  \begin{tabular}{lccc}
    \toprule
    \textbf{Model} & \textbf{Params} & \textbf{Quantization} & \textbf{Disk size} \\
    \midrule
    Qwen3-VL:30b          & 30B & 4-bit & 19\,GB \\
    Gemma3:27b            & 27B & 4-bit & 17\,GB \\
    LLaMA3.2-Vision:11b   & 11B & 4-bit & 8\,GB \\
    \bottomrule
  \end{tabular}
  \vfill
  \begin{itemize}
    \item All served locally via \textbf{Ollama}, single 24\,GB GPU
    \item Deterministic decoding (\texttt{temperature=0}), fully zero-shot: no fine-tuning, no few-shot examples
  \end{itemize}
\end{frame}

\begin{frame}{Evaluation Protocol}
  \begin{description}
    \item[Phase A] Text-only: do models know the formation \emph{names} at all, with no image?
    \item[Phase B] Basic image understanding: jersey colors, offense side, player counts, QB alignment
    \item[Phase C] Formation classification: the core task, in three prompting strategies
  \end{description}
  \vfill
  {\small Majority-class baseline: always predict TRIO $\Rightarrow$ \textbf{35.0\% accuracy}}
\end{frame}

\begin{frame}{Phase C: Three Prompting Strategies}
  \centering
  \begin{columns}[T,onlytextwidth]
    \begin{column}{0.33\textwidth}
      \begin{block}{C1}
        \scriptsize Baseline: image + class names only
      \end{block}
    \end{column}
    \begin{column}{0.33\textwidth}
      \begin{block}{C2}
        \scriptsize + natural-language description per class
      \end{block}
    \end{column}
    \begin{column}{0.33\textwidth}
      \begin{block}{C3}
        \scriptsize + step-by-step decision tree
      \end{block}
    \end{column}
  \end{columns}
  \vfill
 
\end{frame}

\begin{frame}{Pipeline Overview}
  \centering
  \begin{tikzpicture}[
      scale=0.95, every node/.style={scale=0.95},
      node distance=0.5cm and 0.5cm,
      block/.style={rectangle, draw, rounded corners, fill=lite,
                    text width=3.1cm, align=center, minimum height=1.3cm, font=\small},
      arrow/.style={thick, -{Stealth}}
    ]
    \node[block] (clip)   {AFL broadcast clip};
    \node[block, right=of clip]  (ffmpeg) {\texttt{ffmpeg} extraction};
    \node[block, right=of ffmpeg] (image)  {Pre-snap JPEG};

    \node[block, below=of clip]  (vlm)   {Local VLM via Ollama};
    \node[block, right=of vlm]   (prompt) {Phase A/B/C prompting};
    \node[block, right=of prompt] (label) {Regex label extraction};

    \draw[arrow] (clip) -- (ffmpeg);
    \draw[arrow] (ffmpeg) -- (image);
    \draw[arrow] (image.south) -- ++(0,-0.4) -| (vlm.north);
    \draw[arrow] (vlm) -- (prompt);
    \draw[arrow] (prompt) -- (label);

    \node[block, fill=head, text=white, below=of prompt] (metrics) {Accuracy};
    \draw[arrow] (label.south) -- ++(0,-0.4) -| (metrics.north);
  \end{tikzpicture}
\end{frame}

% =====================================================================
\section{Results}
% =====================================================================

\begin{frame}{Phase A: Inconsistencies, Before Any Image (1/2)}
  \centering
  \small
  \renewcommand{\arraystretch}{1.25}
  \begin{tabular}{llccc}
    \toprule
    \textbf{Q} & \textbf{Asked} & \textbf{Qwen3-VL} & \textbf{Gemma3} & \textbf{LLaMA3.2} \\
    \midrule
    A1 & Describe each formation & Empty & Empty & Empty and Trey\\ 
    A2 & No RBs?        & \checkmark & \checkmark & \checkmark \\
    A3 & TEs or FBs?   & \texttimes & \checkmark & \texttimes \\
    A4 & 3 true WRs?    & \texttimes & \texttimes & \texttimes \\
    A5 & 2 WRs each side?        & \checkmark & \checkmark & \texttimes \\
    A6 & TRIO vs. TREY? 			 & \texttimes & \texttimes & \texttimes \\
    A7 & DOUBLES vs. TANK?		 & \texttimes & \texttimes & \texttimes \\
    \bottomrule
  \end{tabular}
  \note{- Walk the table row by row:\\
  	
  	- A2 (name EMPTY): all three correct.\\
  	
  	- A3 (which formations use a TE/FB): only Gemma correct.\\
  	
  	- A4 (name the 3-WR-one-side formation, i.e. TRIO): all three wrong. all had trey\\
  	
  	- A5 (name DOUBLES): Qwen and Gemma correct, LLaMA wrong.\\
  	
  	- A6 (TRIO vs. TREY, explain the difference): all three wrong.\\
  	
  	- A7 (DOUBLES vs. TANK, explain the difference): all three wrong.}
\end{frame}

\begin{frame}{Phase A: Inconsistencies, Before Any Image (2/2)}
  \begin{alertblock}{Finding}
    EMPTY is standard football terminology: all three get it right. DOUBLES, TRIO, TREY, and TANK are not standardized: models struggle to consistently describe or distinguish them
  \end{alertblock}
\end{frame}

\begin{frame}{Phase B: Basic Image Understanding}
  \begin{itemize}
    \item Sanity checks: jersey colors, offense side, player counts, QB positioning
    \item Are models even reliably reading the image (player counts, empty responses)?
    \item \textbf{Swap analysis}: of all offensive-jersey-color errors, how many are the model naming the \emph{defense's} true color instead?
    \item Does giving the offense side explicitly fix this?
  \end{itemize}
\end{frame}

\begin{frame}{Phase B: Results (1/2)}
  \centering
  \tiny
  \renewcommand{\arraystretch}{1.15}
  \begin{tabular}{lccccc}
    \toprule
    \textbf{Model} & \textbf{Off.\ Side} & \textbf{Yard Line} & \textbf{QB Pos.} & \textbf{WR Count} & \textbf{RB Visible} \\
    \midrule
    Qwen3-VL:30b        & 50\% & 25\% & 18\% &  0\% & 45\% \\
    Gemma3:27b          & 45\% & 28\% & 20\% & 25\% & 88\% \\
    LLaMA3.2-Vision:11b & 50\% & 32\% & 20\% & 15\% & 12\% \\
    \bottomrule
  \end{tabular}

  \vspace{6pt}
  \begin{tabular}{lccc|ccc}
    \toprule
     & \multicolumn{3}{c|}{\textbf{B: side inferred}} & \multicolumn{3}{c}{\textbf{B2: side given}} \\
    \textbf{Model} & Correct & Swapped & Other wrong & Correct & Swapped & Other wrong \\
    \midrule
    Qwen3-VL:30b   & 15 & 11 & 14 & 11 &  5 & 24 \\
    Gemma3:27b     & 11 & 19 & 10 & 10 & 20 & 10 \\
    LLaMA3.2:11b   & 15 & 21 &  4 & 17 & 19 &  4 \\
    \bottomrule
  \end{tabular}
\end{frame}

\begin{frame}{Phase B: Results (2/2)}
  \begin{alertblock}{Finding}
    Basic image reading is inconsistent, and often not really ``reading'' at all.
  \end{alertblock}
\end{frame}

\begin{frame}[fragile]{Code Walkthrough}
  \small
  \begin{itemize}
    \item Live in \texttt{evaluation.ipynb}: Ollama API call and prompts C1 $\rightarrow$ C3
  \end{itemize}
\end{frame}

\begin{frame}{Phase C: Formation Classification Accuracy}
  \centering
  \renewcommand{\arraystretch}{1.3}
  \footnotesize
  \begin{tabular}{lccc}
    \toprule
    \textbf{Model} & \textbf{C1} & \textbf{C2} & \textbf{C3}  \\
    \midrule
    Qwen3-VL:30b        & 10.0\%$\pm$9  & 30.0\%$\pm$14 & 27.5\%$\pm$14 \\
    Gemma3:27b          & 15.0\%$\pm$11 & \textbf{37.5\%$\pm$15} & 12.5\%$\pm$10  \\
    LLaMA3.2-Vision:11b & 12.5\%$\pm$10 & 12.5\%$\pm$10 & 20.0\%$\pm$12 \\
    \midrule
    \textbf{Majority baseline} & \multicolumn{3}{c}{\textbf{35.0\%}} \\
    \bottomrule
  \end{tabular}
  \vfill
  \small $\pm$ = 95\% CI (binomial approx., $n$=40 per cell) -- most pairwise differences between phases and models overlap within noise.
\end{frame}

\begin{frame}{Phase C: Models Default, They Don't Classify}
  \centering
  \renewcommand{\arraystretch}{1.3}
  \footnotesize
\begin{tabular}{lccc}
	\toprule
	\textbf{Formation} & \textbf{Ground truth} & \textbf{Predicted} & \textbf{Ratio} \\
	\midrule
	TRIO    & 35.0\% & 15.8\% & 0.5$\times$ \\
	DOUBLES & 27.5\% & 10.0\% & 0.4$\times$ \\
	EMPTY   & 12.5\% &  1.1\% & 0.1$\times$ \\
	TREY    & 12.5\% & 61.9\% & \textbf{5.0$\times$} \\
	TANK    & 12.5\% & 11.1\% & 0.9$\times$ \\
	\bottomrule
\end{tabular}	
  \vfill
  \small Across all 360 Phase C predictions (3 models $\times$ 3 phases $\times$ 40 images): TREY is guessed 5$\times$ more often than it's actually correct, while EMPTY is almost never guessed at all - a favourite default, not a per-image signal.
\end{frame}


% =====================================================================
\section{Discussion \& Conclusion}
% =====================================================================

\begin{frame}{Why It's Not Working (1/2): Terminology}
  \begin{itemize}
    \item Football formation terminology isn't standardized
    \item Evidence: asked directly, with no image at all (Phase A), all three models inconsistent when answering questions about the formations
  \end{itemize}
\end{frame}

\begin{frame}{Why It's Not Working (2/2): Perception}
  \begin{itemize}
    \item Even the basic ``read the image'' step mostly fails, before formation logic can even start
    \item Identifying which side is even the offense: 45-50\%
    \item Wide-receiver count: 0\% (Qwen), 15\% (LLaMA), 25\% (Gemma: comes from reciting ``3'' on every single image, not counting)
    \item Running back in the backfield: LLaMA answers ``No'', Gemma always ``Yes''
  \end{itemize}
\end{frame}

\begin{frame}{Limitations}
  \footnotesize
  \begin{itemize}
    \item Small test set: 40 images, minority classes have only 5 samples each
    \item Single competition (AFL): may not generalize to other leagues/cameras
    \item All models run at 4-bit quantization, constrained by 24\,GB VRAM
    \item Purely zero-shot: no fine-tuning, no in-context examples
  \end{itemize}
\end{frame}

\begin{frame}{Conclusion RQ1 \& RQ2}
  \small
  \begin{description}
    \item[RQ1] No, the gains that beat baseline (37.5\% vs.\ 35.0\%) are fully explained by label collapse onto one or two defaults, not genuine per-image classification, and are within each other's 95\% CIs. 
    \item[RQ2] No model shows real multi-class discrimination. Peak accuracy doesn't track parameter count (30B and 27B tie at 37.5\%, 11B close behind).
  \end{description}
    \note{
  	\begin{itemize}
  		\item RQ1: can user-grade open-source VLMs identify formations zero-shot?
  		\item RQ2: how do the three models (Qwen-VL, Gemma, LLaMA) compare?
  	\end{itemize}
  }
\end{frame}

\begin{frame}{Conclusion RQ3}
  \small
  \begin{description}
    \item[RQ3] Two compounding failure modes:
  \end{description}
  \begin{itemize}
    \item Terminology isn't standardized
    \item No actual reading of the image
  \end{itemize}
    \note{
  	\begin{itemize}
  		\item RQ3: what are the main failure modes on amateur-quality footage?
  	\end{itemize}
  }
\end{frame}

\begin{frame}{Future Work}
  \begin{itemize}
    \item Few-shot prompting with corrected examples
    \item Fine-tuning on a larger, multi-league dataset
  \end{itemize}
  \vfill
  \begin{block}{Master's thesis}
    Computer-vision pipeline: player detection and bird's-eye-view transform of the field, then a \textbf{rule-based} classifier assigns the formation.
  \end{block}
\end{frame}

\begin{frame}[plain]
  \centering
  {\Large Thank you}\\[0.3cm]
  Questions?
\end{frame}

\end{document}
